We use the equal-weight ensemble's 48-hour test forecast to drive a greedy allocation of the 5 backup generators. Given each generator's \textbf{1{,}000-household cap}, the marginal benefit of placing one in county $c$ at decision step $s$ is the integral of the predicted outage in $c$ over the 48-hour horizon, capped both at 1{,}000 (per generator) and at the county's tracked household maximum:

\[
M_s(c) = \sum_{h=1}^{48} \min\!\bigl( \hat y_{c,t+h} - {\rm consumed}_{s-1}(c,h),\ 1000,\ {\rm tracked\_max}(c) \bigr).
\]

We greedily pick the county maximizing $M_s$, deduct the served capacity, and repeat for $s{=}1,\dots,5$. Multiple generators \emph{may} go to the same county if its 48h served-household-hours still exceed 1{,}000 after the previous draw. Our recommendation is:

\begin{quote}
\verb|[ pending — run scripts.finalize ]|
\end{quote}

\paragraph{Rationale.}
[ rationale table pending ]

\paragraph{Limitations.}
(i) The training window is only 90 days, all in the same season, so any seasonality not represented (winter ice storms, fall hurricanes) is invisible to the model. (ii) No future-weather information is available at test time, which fundamentally limits 48-hour predictability of fast-moving convective events. (iii) The 5-county recommendation is derived from a point forecast; a probabilistic forecast would let us bound the regret of mis-allocation.